% !TeX program = xelatex
% 数译（mathtranslations.org）数学中文译本排版模板（讲义 ctexart 版）
% Keith Conrad: Tensor Products II
%
\documentclass[UTF8,12pt,fontset=none]{ctexart}

% ---------- 页面、字体与数学 ----------
\usepackage[paper=a4paper,left=2.5cm,right=2.5cm,top=2.8cm,bottom=2.8cm]{geometry}
\usepackage{fontspec}
\usepackage{xeCJK}
\usepackage{amsmath,amsthm,amssymb,mathrsfs,mathtools}
\usepackage{graphicx,tikz,tikz-cd,titlesec}
\usepackage{array,booktabs,longtable}
\usepackage{etoolbox}
\usepackage[dvipsnames]{xcolor}
\usepackage[
    colorlinks=true,
    bookmarksnumbered=true,
    linkcolor=MidnightBlue,
    citecolor=BrickRed,
    urlcolor=MidnightBlue,
    linktoc=all
]{hyperref}

% 允许版心稍微拉伸以避免孤立过宽行
\setlength{\emergencystretch}{3em}

% 英文：Computer Modern Unicode Serif
\setmainfont{cmunrm.otf}[
  BoldFont=cmunbx.otf,
  ItalicFont=cmunti.otf,
  BoldItalicFont=cmunbi.otf
]

% 中文：正文宋体；专用术语楷体；定理类环境仿宋
\setCJKmainfont{FandolSong-Regular.otf}[BoldFont=FandolSong-Bold.otf]
\newCJKfontfamily\songfont{FandolSong-Regular.otf}[BoldFont=FandolSong-Bold.otf]
\newCJKfontfamily\termfont{FandolKai-Regular.otf}
\newCJKfontfamily\theoremfont[AutoFakeBold=2]{FandolFang-Regular.otf}
\setCJKsansfont{FandolHei-Regular.otf}[BoldFont=FandolHei-Bold.otf]
\setCJKmonofont{FandolFang-Regular.otf}

\pagestyle{plain}
\setlength{\parindent}{2em}
\setlength{\parskip}{0pt}
\linespread{1.16}

% ---------- 译本元信息 ----------
\newcommand{\BookTitleCN}{张量积 II}
\newcommand{\BookTitleEN}{Tensor Products II}
\newcommand{\OriginalAuthor}{Keith Conrad}
\newcommand{\OriginalEdition}{Expository Notes}
\newcommand{\OriginalPublisher}{University of Connecticut}
\newcommand{\OriginalYear}{2024}
\newcommand{\Translator}{数译}
\newcommand{\ModelUsed}{GLM-5 / Codex}
\newcommand{\TranslationDate}{\today}

% ================================================================
%                         术语系统
% ================================================================
\newcommand{\termcn}[1]{{\termfont #1}}

\makeatletter
\newcommand{\MTTermList}{}
\newcommand{\MTRegisterTerm}[3]{%
  \g@addto@macro\MTTermList{\MTTermRow{#1}{#2}{#3}}%
}
\newcommand{\MTTermRow}[3]{%
  {\termfont #1} & #2 & \hyperref[#3]{\pageref*{#3}} \\
}

\newcommand{\newterm}[3]{%
  \@ifundefined{MT@seen@#1}{%
    \expandafter\gdef\csname MT@seen@#1\endcsname{}%
    \phantomsection
    \label{term:#1}%
    \protected@write\@auxout{}{%
      \string\MTRegisterTerm{\unexpanded{#2}}{\unexpanded{#3}}{term:#1}%
    }%
  }{}%
  {\termfont #2}（#3）%
}
\makeatother

\newcommand{\printterminology}{%
  \clearpage
  \phantomsection
  \section*{术语索引}
  \addcontentsline{toc}{section}{术语索引}

  \vspace{0.5cm}
  \begingroup
  \renewcommand{\arraystretch}{1.28}
  \begin{longtable}{>{\raggedright\arraybackslash}p{0.28\textwidth}>{\raggedright\arraybackslash}p{0.49\textwidth}>{\centering\arraybackslash}p{0.12\textwidth}}
    \toprule
    \textbf{中文术语} & \textbf{英文术语} & \textbf{页码} \\
    \midrule
    \endfirsthead
    \toprule
    \textbf{中文术语} & \textbf{英文术语} & \textbf{页码} \\
    \midrule
    \endhead
    \MTTermList
    \bottomrule
  \end{longtable}
  \endgroup
}

% ================================================================
%                       定理类环境
% ================================================================
\newtheoremstyle{mtfang}%
  {6pt}{6pt}%
  {\theoremfont}%
  {0pt}%
  {\songfont\bfseries}%
  {：}%
  {0.5em}%
  {}%

\newtheoremstyle{mtsong}%
  {6pt}{6pt}%
  {\songfont}%
  {0pt}%
  {\songfont\bfseries}%
  {：}%
  {0.5em}%
  {}%

% 讲义中公式随节编号
\numberwithin{equation}{section}

% 所有定理类环境共用节级计数器，与原书完全一致
\theoremstyle{mtfang}
\newtheorem{theorem}{定理}[section]
\newtheorem{lemma}[theorem]{引理}
\newtheorem{proposition}[theorem]{命题}
\newtheorem{corollary}[theorem]{推论}
\newtheorem{definition}[theorem]{定义}
\newtheorem{remark}[theorem]{注}

\theoremstyle{mtsong}
\newtheorem{example}[theorem]{例}
\newtheorem{problem}[theorem]{问题}
\newtheorem{question}[theorem]{问题}

\newcommand{\envenddiamond}{%
  \ifhmode
    \unskip\nobreak\hfill\ensuremath{\diamond}%
  \else
    \noindent\hfill\ensuremath{\diamond}%
  \fi
  \par
}
\AtEndEnvironment{example}{\envenddiamond}
\AtEndEnvironment{remark}{\envenddiamond}

% proof 样式
\makeatletter
\renewcommand{\proofname}{证明}
\renewenvironment{proof}[1][\proofname]{\par
  \pushQED{\qed}%
  \normalfont\songfont
  \topsep6\p@\@plus6\p@\relax
  \trivlist
  \item[\hskip\labelsep\normalfont\songfont\bfseries #1：]\ignorespaces
}{%
  \popQED\endtrivlist\@endpefalse
}
\makeatother
\renewcommand{\qedsymbol}{\ensuremath{\square}}

% ================================================================
%                      讲义节标题格式
% ================================================================
\titleformat{\section}[block]
  {\normalfont\Large\bfseries}
  {\S\,\thesection\quad}
  {0pt}
  {}
\titlespacing*{\section}{0pt}{2ex plus 0.5ex minus .2ex}{1.2ex plus .2ex}

\ctexset{
  subsection={
    format=\large\bfseries,
    name={,.\quad},
    number=\arabic{section}.\arabic{subsection}
  }
}

% 常用数学记号
\newcommand{\R}{\mathbf{R}}
\newcommand{\C}{\mathbf{C}}
\newcommand{\Z}{\mathbf{Z}}
\newcommand{\Q}{\mathbf{Q}}
\newcommand{\N}{\mathbf{N}}
\newcommand{\F}{\mathbf{F}}
\newcommand{\Hom}{\operatorname{Hom}}
\newcommand{\End}{\operatorname{End}}
\newcommand{\Aut}{\operatorname{Aut}}
\newcommand{\im}{\operatorname{im}}
\newcommand{\id}{\operatorname{id}}
\newcommand{\coker}{\operatorname{coker}}
\newcommand{\Tor}{\operatorname{Tor}}
\newcommand{\Ext}{\operatorname{Ext}}
\newcommand{\ann}{\operatorname{ann}}
\newcommand{\Tr}{\operatorname{Tr}}

% 参考文献环境
\renewcommand{\refname}{参考文献}

\begin{document}

% ================================================================
%                             封面
% ================================================================
\begin{titlepage}
    \thispagestyle{empty}
    \flushleft

    \rule{\textwidth}{1.2pt} \\[0.1cm]
    \rule{\textwidth}{0.8pt} \\[1.5cm]

    \hspace{2cm} {\fontsize{42pt}{52pt}\selectfont \bfseries \BookTitleCN} \\[0.8cm]
    \hspace{2cm} {\fontsize{24pt}{28pt}\selectfont \textit{\BookTitleEN}} \\[1.5cm]

    \rule{\textwidth}{0.6pt} \\[2cm]

    \hspace{2cm} {\Huge \textsc{\OriginalAuthor}} \\[5cm]

    \vfill

    \noindent
    \begin{minipage}[b]{0.46\textwidth}
        \raggedright
        \includegraphics[width=5.7cm]{logo.pdf}
    \end{minipage}
    \hfill
    \begin{minipage}[b]{0.50\textwidth}
        \raggedleft
        {\small \OriginalEdition · \OriginalPublisher · \OriginalYear} \\[0.25cm]
        {\normalsize \Translator\ 翻译及重排} \\[0.25cm]
        {\small Model: \ModelUsed} \\
        {\small Updated: \TranslationDate}
    \end{minipage}

    \vspace*{1cm}
\end{titlepage}

% ================================================================
%                       译者说明与目录
% ================================================================
\pagenumbering{roman}
\setcounter{page}{1}

\phantomsection
\section*{译者说明}
\addcontentsline{toc}{section}{译者说明}

本讲义是康涅狄格大学（University of Connecticut）Keith Conrad 教授所撰写的高等代数著名讲义《张量积 II》（\textit{Tensor Products II}）的中译本，是前篇《张量积》（Part I）的自然深化与进阶.

在第一篇确立了模与向量空间张量积的基本定义、普遍映射性质及计算范例的基础上，本篇讲义全面转向更深层的代数结构：
\begin{enumerate}
    \item \textbf{映射的张量积}：系统考察线性映射的张量积 $\varphi \otimes \psi$，其在矩阵层面上导出的 Kronecker 积，以及由其诱导的算子特征值乘法性质；
    \item \textbf{平坦模与正合性}：引入同调代数与交换代数的核心概念——平坦模（Flat Modules），深入剖析张量积函子的右正合性与非左正合现象；
    \item \textbf{基变换下的态射与代数结构}：研究基变换如何保持态射结构，探讨向量空间张量积的独特性质、张量缩并（Contraction）运算；
    \item \textbf{代数的张量积与模的张量代数}：详尽阐述交换环上代数的张量积运算（$R$-代数张量积），以及从任意模出发构造非交换自由分次代数——张量代数 $T(M)$ 的万有性质.
\end{enumerate}

排版规范严格继承第一篇：定理类环境共享统一连续计数器，所有交换图均由 \texttt{tikz-cd} 矢量重绘，行内公式与跨章节交叉引用完整健全，并在文末提供中英对照术语索引.

\vspace{0.5cm}
\hfill \Translator

\hfill \TranslationDate

\newpage
\begingroup
\hypersetup{linkcolor=black}
\tableofcontents
\endgroup
\clearpage

\pagenumbering{arabic}
\setcounter{page}{1}

% sections2/sec1_intro.tex
\section{引言}
\label{sec:intro}

继续我们对张量积的研究，我们将探讨如何将两个线性映射 $M \to M'$ 与 $N \to N'$ 组合为一个线性映射 $M \otimes_R N \to M' \otimes_R N'$.由此将自然引出\newterm{flat-module}{平坦模}{flat module}的概念以及基变换之间的线性映射.随后，我们将考察向量空间张量积的独特性质（包括张量缩并）、$R$-代数的张量积，并最终讨论 $R$-模的张量代数.

% sections2/sec2_linear_maps.tex
\section{线性映射的张量积}
\label{sec:linear-maps}

若 $M \xrightarrow{\varphi} M'$ 与 $N \xrightarrow{\psi} N'$ 为 $R$-线性映射，则在直和之间存在由 $(\varphi \oplus \psi)(m, n) = (\varphi(m), \psi(n))$ 给出的线性映射 $M \oplus N \xrightarrow{\varphi \oplus \psi} M' \oplus N'$.我们希望在张量积之间同样定义一个线性映射 $M \otimes_R N \to M' \otimes_R N'$，使得简单张量满足
\[
  m \otimes n \mapsto \varphi(m) \otimes \psi(n).
\]

从映射 $M \times N \to M' \otimes_R N'$（定义为 $(m, n) \mapsto \varphi(m) \otimes \psi(n)$）出发，容易验证它是 $R$-双线性的.因此由张量积的普遍映射性质，存在唯一的 $R$-线性映射 $M \otimes_R N \xrightarrow{\varphi \otimes \psi} M' \otimes_R N'$，使得 $(\varphi \otimes \psi)(m \otimes n) = \varphi(m) \otimes \psi(n)$，更一般地作用在有限和张量上为：
\[
  (\varphi \otimes \psi)(m_1 \otimes n_1 + \dots + m_k \otimes n_k) = \varphi(m_1) \otimes \psi(n_1) + \dots + \varphi(m_k) \otimes \psi(n_k).
\]
我们称 $\varphi \otimes \psi$ 为 $\varphi$ 与 $\psi$ 的\newterm{tensor-prod-maps}{张量积映射}{tensor product of linear maps}.需要特别辨析的是：$\varphi \otimes \psi$ 并不是某个模张量积中的简单张量，它仅仅代表作用在 $M \otimes_R N$ 上的全新线性映射记号.

特别地，当取其中一个映射为恒等映射时，映射 $N \otimes_R M \xrightarrow{1 \otimes \varphi} N \otimes_R M'$ 或 $M \otimes_R N \xrightarrow{\varphi \otimes 1} M' \otimes_R N$ 称为\termcn{与 $N$ 作张量积}（tensoring with $N$）.这一运算在第~\ref{sec:linear-base-ext}~节讨论基变换时具有基础性的核心地位.

\begin{example}\label{ex:inclusion-tensor}
考虑包含映射 $a\Z \xrightarrow{i} \Z$ 与 $\Z/b\Z$ 作张量积：$a\Z \otimes_\Z \Z/b\Z \xrightarrow{i \otimes 1} \Z \otimes_\Z \Z/b\Z$.在 $\Z \otimes_\Z \Z/b\Z \cong \Z/b\Z$ 中，$(i \otimes 1)(ax \otimes y \bmod b) = ax \otimes (y \bmod b) = 1 \otimes (axy \bmod b)$.若令 $d = \gcd(a, b)$，则 $i \otimes 1$ 的像由 $1 \otimes (d \bmod b)$ 生成.因此当 $\gcd(a, b) = 1$ 时 $i \otimes 1$ 是满射，而当 $b \mid a$ 时其像为 $0$.特别地，当 $a \geq 2$ 且 $\gcd(a, b) = 1$ 时，$i$ 本身不是满射，但张量积映射 $i \otimes 1$ 却变成了满射！
\end{example}

\begin{example}[Kronecker 积]\label{ex:kronecker-prod}
设 $A = \begin{pmatrix} a & b \\ c & d \end{pmatrix}$ 与 $A' = \begin{pmatrix} a' & b' \\ c' & d' \end{pmatrix}$ 属于 $\mathrm{M}_2(R)$.它们均诱导线性映射 $R^2 \to R^2$.因此 $A \otimes A'$ 是 $(R^2)^{\otimes 2}$ 上的线性变换.在标准基 $\{e_1 \otimes e_1, e_1 \otimes e_2, e_2 \otimes e_1, e_2 \otimes e_2\}$ 下，计算 $A \otimes A'$ 的矩阵表示：
\begin{align*}
  (A \otimes A')(e_1 \otimes e_1) &= Ae_1 \otimes A'e_1 = (ae_1 + ce_2) \otimes (a'e_1 + c'e_2) \\
  &= aa' e_1 \otimes e_1 + ac' e_1 \otimes e_2 + ca' e_2 \otimes e_1 + cc' e_2 \otimes e_2,
\end{align*}
类似地算出其余各列，得到 $4 \times 4$ 分块矩阵：
\[
  A \otimes A' = \begin{pmatrix}
    aa' & ab' & ba' & bb' \\
    ac' & ad' & bc' & bd' \\
    ca' & cb' & da' & db' \\
    cc' & cd' & dc' & dd'
  \end{pmatrix} = \begin{pmatrix}
    a A' & b A' \\
    c A' & d A'
  \end{pmatrix}.
\]
其迹满足 $\Tr(A \otimes A') = (\Tr A)(\Tr A')$.

更一般地，对于 $A \in \mathrm{M}_m(R)$ 与 $B \in \mathrm{M}_n(R)$，分块矩阵 $(a_{ij} B)$ 称为 $A$ 与 $B$ 的\newterm{kronecker-product}{Kronecker 积}{Kronecker product}.
\end{example}

\begin{theorem}[特征值乘积定理]\label{thm:eigen-tensor}
设 $K$ 为域，$A \in \mathrm{M}_m(K)$ 与 $B \in \mathrm{M}_n(K)$ 分别具有特征值 $\lambda$ 与 $\mu$.则 $A \otimes I_n + I_m \otimes B$ 具有特征值 $\lambda + \mu$，而 $A \otimes B$ 具有特征值 $\lambda \mu$.
\end{theorem}

\begin{proof}
设 $Av = \lambda v$ 且 $Bw = \mu w$（$v, w \neq 0$）.则在简单张量 $v \otimes w \neq 0$ 上直接计算：
\[
  (A \otimes I_n + I_m \otimes B)(v \otimes w) = Av \otimes w + v \otimes Bw = \lambda(v \otimes w) + \mu(v \otimes w) = (\lambda + \mu)(v \otimes w),
\]
且 $(A \otimes B)(v \otimes w) = Av \otimes Bw = \lambda \mu(v \otimes w)$.
\end{proof}

\begin{example}\label{ex:poly-roots-tensor}
$\sqrt{2}$ 与 $\sqrt{3}$ 分别为 $A = \begin{pmatrix} 0 & 2 \\ 1 & 0 \end{pmatrix}$ 与 $B = \begin{pmatrix} 0 & 3 \\ 1 & 0 \end{pmatrix}$ 的特征值.由定理~\ref{thm:eigen-tensor}，矩阵 $A \otimes I_2 + I_2 \otimes B$ 必定具有特征值 $\sqrt{2} + \sqrt{3}$.计算该 $4 \times 4$ 矩阵的特征多项式即得 $T^4 - 10T^2 + 1$，由此纯代数地给出了以 $\sqrt{2} + \sqrt{3}$ 为根的有理系数多项式！
\end{example}

\begin{theorem}\label{thm:hom-tensor-hom}
存在自然的 $R$-线性映射
\[
  \Hom_R(M, M') \otimes_R \Hom_R(N, N') \to \Hom_R(M \otimes_R N, M' \otimes_R N'),
\]
将简单张量 $\varphi \otimes \psi$ 映为线性映射 $\varphi \otimes \psi$.当 $M, M', N, N'$ 均为有限生成自由模时，该映射为同构.
\end{theorem}

\begin{proof}
从双线性映射 $(\varphi, \psi) \mapsto [\text{映射 } m \otimes n \mapsto \varphi(m) \otimes \psi(n)]$ 出发，由普遍映射性质即得正向线性映射.当模具有有限基时，初等矩阵基底的一一对应立即可知其将基映为基，故为同构.
\end{proof}

\begin{theorem}[函子性与复合律]\label{thm:functoriality}
对于 $R$-模 $M$ 与 $N$，恒有 $\id_M \otimes \id_N = \id_{M \otimes_R N}$.对任意可复合的线性映射序列：
\[
  (\varphi' \otimes \psi') \circ (\varphi \otimes \psi) = (\varphi' \circ \varphi) \otimes (\psi' \circ \psi).
\]
\end{theorem}

\begin{example}[张量积行列式公式]\label{ex:det-tensor-maps}
设 $M, N$ 为有限自由 $R$-模，秩分别为 $k$ 与 $\ell$.对于线性算子 $\varphi \colon M \to M$ 与 $\psi \colon N \to N$，利用复合律可将 $\varphi \otimes \psi$ 分解为
\[
  \varphi \otimes \psi = (\varphi \otimes \id_N) \circ (\id_M \otimes \psi).
\]
选取适当基底可知 $\varphi \otimes \id_N$ 对应于对角线上由 $\ell$ 个矩阵块 $[\varphi]$ 构成的分块对角阵，故其行列式为 $(\det \varphi)^\ell$.因此：
\begin{equation}
  \det(\varphi \otimes \psi) = (\det \varphi)^\ell (\det \psi)^k. \label{eq:2.1}
\end{equation}
特别地，在例~\ref{ex:kronecker-prod} 中 $\det(A \otimes A') = (\det A)^2 (\det A')^2$.
\end{example}

\begin{corollary}\label{cor:tensor-power-det}
设 $M$ 为秩 $k \geq 1$ 的自由模，$\varphi \colon M \to M$ 为线性算子.对所有 $i \geq 1$：
\[
  \det(\varphi^{\otimes i}) = (\det \varphi)^{i k^{i-1}}.
\]
\end{corollary}

\begin{remark}\label{rem:kronecker-zehfuss}
公式~\eqref{eq:2.1} 在矩阵层面上即为 $\det(A \otimes B) = (\det A)^\ell (\det B)^k$.虽然 $A \otimes B$ 通称为 Kronecker 积，但行列式公式在历史上最早是由 Zehfuss 给出的，参见 Henderson 与 Pukelsheim \cite{henderson-pukelsheim}.
\end{remark}

\begin{theorem}[同构与满射的保持性]\label{thm:iso-surj-preserved}
\begin{enumerate}
  \item 若 $\varphi \colon M \to M'$ 与 $\psi \colon N \to N'$ 均为同构，则 $\varphi \otimes \psi$ 为同构（其逆为 $\varphi^{-1} \otimes \psi^{-1}$）.
  \item 若 $\varphi$ 与 $\psi$ 均为满射，则 $\varphi \otimes \psi$ 为满射.
\end{enumerate}
\end{theorem}

\textbf{极其重要的核心反例：张量积通常不保持单射！}即使 $\varphi$ 与 $\psi$ 均单，$\varphi \otimes \psi$ 也完全可能不是单射.

\begin{example}\label{ex:non-injective-tensor}
设 $R = \Z$，考虑数乘映射 $\alpha \colon \Z/p\Z \to \Z/p^2\Z$（$\alpha(x) = px$）.显然 $\alpha$ 是单射.然而将其与 $\Z/p\Z$ 作张量积：
\[
  1 \otimes \alpha \colon \Z/p\Z \otimes_\Z \Z/p\Z \to \Z/p\Z \otimes_\Z \Z/p^2\Z, \quad a \otimes x \mapsto a \otimes px = pa \otimes x = 0.
\]
映射 $1 \otimes \alpha$ 恒等于零！由于定义域 $\Z/p\Z \otimes_\Z \Z/p\Z \cong \Z/p\Z \neq 0$，该映射绝不是单射.
\end{example}

这一根本特性说明：若 $M \subset M'$ 且 $N \subset N'$，\textbf{子模的张量积 $M \otimes_R N$ 通常不能直接视作 $M' \otimes_R N'$ 的子模}，因为自然包含映射的张量积可能出现灾难性的坍缩！

\begin{example}\label{ex:tensor-square-not-injective}
设 $R = A[X, Y]$（$A$ 为交换环），考虑理想 $I = (X, Y)$.包含映射 $i \colon I \hookrightarrow R$ 是单射，但其张量平方 $i^{\otimes 2} \colon I^{\otimes 2} \to R^{\otimes 2}$ 却\textbf{不是单射}.

在 $R^{\otimes 2}$ 中，$X \otimes Y = XY(1 \otimes 1) = Y \otimes X$，故 $X \otimes Y - Y \otimes X = 0$.但在 $I^{\otimes 2}$ 中，$X \otimes Y \neq Y \otimes X$！因为映射 $(f, g) \mapsto f_X(0, 0)g_Y(0, 0)$ 是良定义的双线性映射并诱导线性映射 $I^{\otimes 2} \to A$，它把 $X \otimes Y$ 映为 $1$，而把 $Y \otimes X$ 映为 $0$.因此非零张量 $X \otimes Y - Y \otimes X$ 在包含映射下落入了核中.
\end{example}

\begin{remark}\label{rem:torsion-free-tensor-torsion}
例~\ref{ex:tensor-square-not-injective} 中的理想 $I = (X, Y)$ 在整环上是无扭模，但 $I^{\otimes 2}$ 却带有非零扭元：$X(X \otimes Y - Y \otimes X) = 0$！因此\textbf{无扭模的张量积甚至未必是无扭的}.
\end{remark}

\begin{theorem}\label{thm:direct-summand-injective}
设 $\varphi \colon M \to N$ 为单射，且 $\varphi(M)$ 是 $N$ 的直和项（即序列分裂）.则对任意 $k \geq 0$，$\varphi^{\otimes k} \colon M^{\otimes k} \to N^{\otimes k}$ 均为单射，且其像亦为直和项.
\end{theorem}

\begin{theorem}[满同态张量积的核]\label{thm:kernel-surjective-tensor}
设 $\varphi \colon M \to M'$ 与 $\psi \colon N \to N'$ 均为满线性同态.则张量积映射 $\varphi \otimes \psi \colon M \otimes_R N \to M' \otimes_R N'$ 的核由所有满足 $\varphi(m) = 0$ 或 $\psi(n) = 0$ 的简单张量 $m \otimes n$ 生成：
\[
  \ker(\varphi \otimes \psi) = (i \otimes 1)((\ker \varphi) \otimes_R N) + (1 \otimes j)(M \otimes_R (\ker \psi)).
\]
\end{theorem}

\begin{corollary}\label{cor:quotient-ideal-tensor}
设 $I$ 为 $R$ 的理想，则
\[
  (M \otimes_R N) / I(M \otimes_R N) \cong (M/IM) \otimes_R (N/IN).
\]
\end{corollary}

\begin{corollary}\label{cor:poly-quotient-base-ext}
设 $f \colon R \to S$ 为环同态，$I$ 为多项式环 $R[X_1, \dots, X_n]$ 的理想.则作为 $S$-模（以及 $S$-代数）有
\[
  S \otimes_R \left(R[X_1, \dots, X_n] / I\right) \cong S[X_1, \dots, X_n] / (I \cdot S[X_1, \dots, X_n]).
\]
\end{corollary}

% sections2/sec3_flat.tex
\section{平坦模}
\label{sec:flat}

既然线性映射的张量积通常不能保持单射性，我们就必须特别关注并命名那些能够始终保持单射性的 $R$-模 $N$：即只要 $M \xrightarrow{\varphi} M'$ 是单射，由其诱导的映射 $N \otimes_R M \xrightarrow{1 \otimes \varphi} N \otimes_R M'$ 就必定也是单射（注意在 $N$ 上的映射是恒等映射）.

\begin{definition}\label{def:flat-module}
设 $N$ 为 $R$-模.若对所有单射线性映射 $M \xrightarrow{\varphi} M'$，诱导的线性映射
\[
  N \otimes_R M \xrightarrow{1 \otimes \varphi} N \otimes_R M'
\]
恒为单射，则称 $N$ 为\newterm{flat-mod-def}{平坦模}{flat module}.
\end{definition}

平坦模的概念之所以重要，在于自然界中存在极为丰富且核心的平坦模范例：

\begin{theorem}[自由模必平坦]\label{thm:free-is-flat}
任何自由 $R$-模 $F$ 都是平坦模：若线性映射 $\varphi \colon M \to M'$ 是单射，则
\[
  1 \otimes \varphi \colon F \otimes_R M \to F \otimes_R M'
\]
必为单射.
\end{theorem}

\begin{proof}
当 $F = 0$ 时显然.设 $F \neq 0$ 且具有一组基 $\{e_i\}_{i \in I}$.在第一篇中我们已经证明， $F \otimes_R M$ 中的每个元素都可以唯一地表示为 $\sum_{i} e_i \otimes m_i$（其中 $m_i \in M$），对 $F \otimes_R M'$ 亦然.

设 $t \in \ker(1 \otimes \varphi)$，将 $t$ 展开为 $\sum_i e_i \otimes m_i$.则
\[
  0 = (1 \otimes \varphi)(t) = \sum_{i} e_i \otimes \varphi(m_i).
\]
由基展开表示的唯一性，必有每个 $\varphi(m_i) = 0$.因为 $\varphi$ 是单射，故所有 $m_i = 0$，从而 $t = 0$.
\end{proof}

注意：定理~\ref{thm:free-is-flat} 对自由模的基没有任何有限性要求，无限基自由模同样平坦.

\begin{theorem}[分式域必平坦]\label{thm:fraction-field-flat}
设 $R$ 为整环，$K$ 为其分式域.则 $K$ 作为 $R$-模是平坦模.
\end{theorem}

\begin{remark}\label{rem:fraction-not-free}
定理~\ref{thm:fraction-field-flat} 绝非定理~\ref{thm:free-is-flat} 的特例：若 $K$ 作为 $R$-模是自由模，则必然要求 $K = R$（因为 $K$ 中任意两元素均 $R$-线性相关，其自由秩至多为 $1$；若为 $1$ 则 $K \cong R$，从而 $1/a \in R$，要求 $R$ 是域）.因此只要 $R$ 不是域（如 $R = \Z$），其分式域 $K$（如 $\Q$）便给出了一个\textbf{平坦却非自由}的经典范例！
\end{remark}

\begin{proof}[定理~\ref{thm:fraction-field-flat} 的证明]
设 $M \xrightarrow{\varphi} M'$ 为单射 $R$-模同态.回忆第一篇证明过的结论：$K \otimes_R M$ 中的每个张量都是简单张量（通分即可），且简单张量 $x \otimes m = 0$ 当且仅当 $x = 0$ 或 $m \in M_{\mathrm{tor}}$.

设 $(1 \otimes \varphi)(t) = 0$.将 $t$ 写为 $x \otimes m$，则 $0 = (1 \otimes \varphi)(t) = x \otimes \varphi(m)$.因此 $x = 0$ 或者 $\varphi(m) \in M'_{\mathrm{tor}}$.若 $\varphi(m) \in M'_{\mathrm{tor}}$，则存在非零 $r \in R$ 使得 $r\varphi(m) = 0$.由线性性 $\varphi(rm) = 0$；由 $\varphi$ 单射即得 $rm = 0$，故 $m \in M_{\mathrm{tor}}$.因此在任何情况下均有 $t = x \otimes m = 0$，证明了 $1 \otimes \varphi$ 是单射.
\end{proof}

若 $M$ 是 $M'$ 的子模，定理~\ref{thm:fraction-field-flat} 保证了我们可以自然地将 $K \otimes_R M$ 视为 $K \otimes_R M'$ 的子空间.即使 $M$ 或 $M'$ 内部带有扭元（此时 $M \to K \otimes_R M$ 自身可能不是单射），张量积之间的映射 $K \otimes_R M \to K \otimes_R M'$ 却依然保持严格单射！

\begin{remark}\label{rem:localization-flat}
定理~\ref{thm:fraction-field-flat} 可直接推广：对于任意交换环 $R$ 及其乘法子集 $D$，局部化环 $R_D$ 作为 $R$-模恒为平坦模.
\end{remark}

\begin{theorem}\label{thm:flat-ideal-mult}
若 $M$ 是平坦 $R$-模，则对 $R$ 的任意理想 $I$，均有同构
\[
  I \otimes_R M \cong IM, \quad i \otimes m \mapsto im.
\]
\end{theorem}

\begin{proof}
包含映射 $I \hookrightarrow R$ 是单射.由 $M$ 的平坦性，诱导单射 $I \otimes_R M \hookrightarrow R \otimes_R M \cong M$.复合映射将 $i \otimes m$ 映为 $im$，其像恰好是子模 $IM$，故为同构.
\end{proof}

\begin{example}[非平坦模范例]\label{ex:non-flat-torsion}
非零扭阿贝尔群 $A$（例如有限群 $\Z/n\Z$ 或无限群 $\Q/\Z$）作为 $\Z$-模绝不是平坦的：包含映射 $\Z \hookrightarrow \Q$ 是单射，但与其作张量积后得到 $A \otimes_\Z \Z \cong A \to A \otimes_\Z \Q = 0$，单射性彻底破坏！更一般地，整环上的任何非零扭模都不是平坦模.
\end{example}

\begin{theorem}[平坦模与核的精确保持]\label{thm:flat-preserves-kernel}
设 $N$ 为平坦 $R$-模，$\varphi \colon M \to M'$ 为线性映射.则映射
\[
  N \otimes_R M \xrightarrow{1 \otimes \varphi} N \otimes_R M'
\]
的核恰好等于 $N \otimes_R \ker \varphi$（自然地视作 $N \otimes_R M$ 的子模）：
\[
  \ker(1 \otimes \varphi) = (1 \otimes i)(N \otimes_R \ker \varphi).
\]
换言之，\textbf{与平坦模作张量积严格保持任意线性映射的核}.
\end{theorem}

\begin{theorem}\label{thm:flat-tensor-flat}
两个平坦模的张量积依然是平坦模.
\end{theorem}

\begin{proof}
设 $N$ 与 $N'$ 平坦，$\varphi \colon M \to M'$ 为单射.由 $N'$ 平坦，$1 \otimes \varphi \colon N' \otimes_R M \to N' \otimes_R M'$ 为单射；再由 $N$ 平坦，$1 \otimes (1 \otimes \varphi) \colon N \otimes_R (N' \otimes_R M) \to N \otimes_R (N' \otimes_R M')$ 为单射.利用结合律同构，即得 $(N \otimes_R N') \otimes_R M \to (N \otimes_R N') \otimes_R M'$ 为单射.
\end{proof}

\begin{theorem}[双边映射单射判别准则]\label{thm:injective-tensor-both}
设 $\varphi \colon M \to M'$ 与 $\psi \colon N \to N'$ 均为单射.若满足下列条件之一：
\begin{enumerate}
  \item $M$ 与 $N'$ 均为平坦模；
  \item $M'$ 与 $N$ 均为平坦模；
\end{enumerate}
则张量积映射 $\varphi \otimes \psi \colon M \otimes_R N \to M' \otimes_R N'$ 必为单射.
\end{theorem}

\begin{proof}
将 $\varphi \otimes \psi$ 分解为复合：
\[
  M \otimes_R N \xrightarrow{1 \otimes \psi} M \otimes_R N' \xrightarrow{\varphi \otimes 1} M' \otimes_R N'.
\]
若 $M$ 平坦，则 $1 \otimes \psi$ 单；若 $N'$ 平坦，则 $\varphi \otimes 1$ 单.两单射之复合必为单射！对另一种条件对称分解即可.
\end{proof}

\begin{corollary}\label{cor:flat-tensor-powers-injective}
设 $M_1, \dots, M_k$ 与 $N_1, \dots, N_k$ 均为平坦 $R$-模，且 $\varphi_i \colon M_i \to N_i$ 均为单射.则
\[
  \varphi_1 \otimes \dots \otimes \varphi_k \colon M_1 \otimes_R \dots \otimes_R M_k \to N_1 \otimes_R \dots \otimes_R N_k
\]
必为单射.特别地，平坦模之间单射的张量幂 $\varphi^{\otimes k}$ 恒为单射.
\end{corollary}

\begin{corollary}\label{cor:free-independent-tensors}
设 $M$ 为自由 $R$-模，$\{m_1, \dots, m_d\}$ 是 $M$ 中的有限线性无关子集.则对任意 $k \leq d$，$d^k$ 个简单张量
\[
  m_{i_1} \otimes \dots \otimes m_{i_k} \quad (1 \leq i_1, \dots, i_k \leq d)
\]
在张量幂 $M^{\otimes k}$ 中必定是 $R$-线性无关的.
\end{corollary}

% sections2/sec4_linear_base_ext.tex
\section{线性映射的张量积与基变换}
\label{sec:linear-base-ext}

固定环同态 $R \xrightarrow{f} S$.每个 $S$-模均可通过标量限制成为 $R$-模，而每个 $R$-模 $M$ 均具有基变换（标量扩张）$S \otimes_R M$（它是一个 $S$-模）.在第一篇中我们看到，$S \otimes_R M$ 在所有 $S$-模中满足普遍映射性质：从 $M$ 出发到达任意 $S$-模的 $R$-线性映射均可唯一延拓为从 $S \otimes_R M$ 出发的 $S$-线性映射.

在本节中，我们探讨基变换更为重要的另一面：它将两个 $R$-模之间的 $R$-线性映射 $M \xrightarrow{\varphi} M'$ 自然转化为 $S$-模之间的 $S$-线性映射.将映射 $\varphi$ 与 $S$ 作张量积，得到 $R$-线性映射
\[
  S \otimes_R M \xrightarrow{1 \otimes \varphi} S \otimes_R M'.
\]
容易验证该映射事实上是 $S$-线性的：$(1 \otimes \varphi)(st) = s(1 \otimes \varphi)(t)$.为了突显其对 $S$ 的依赖性，我们通常将该基变换后的线性映射记作 $\varphi_S$：
\begin{equation}
  \varphi_S \colon S \otimes_R M \to S \otimes_R M', \quad \varphi_S(s \otimes m) = s \otimes \varphi(m).
\end{equation}
由于 $(\id_M)_S = \id_{S \otimes_R M}$ 且 $(\varphi \circ \varphi')_S = \varphi_S \circ \varphi'_S$，基变换过程具有完整的函子性.

若 $R$-线性映射 $\varphi$ 是同构或满射，则其基变换 $\varphi_S$ 亦保持同构或满射.但若 $\varphi$ 是单射，$\varphi_S$ 却未必是单射（见第一篇及第二节的反例）.

\begin{theorem}\label{thm:invariance-rank-domain}
设 $R$ 为非零交换环.若 $R^m \cong R^n$ 为 $R$-模同构，则 $m = n$；若存在满线性映射 $R^m \twoheadrightarrow R^n$，则 $m \geq n$.
\end{theorem}

\begin{proof}
选取 $R$ 的极大理想 $\mathfrak{m}$.将映射与域 $R/\mathfrak{m}$ 作张量积，即可转化为域上向量空间的线性同构或满射，由向量空间维数理论即得 $m = n$ 或 $m \geq n$.
\end{proof}

\begin{theorem}\label{thm:pid-free-rank}
设 $R$ 为 PID，$M$ 为有限生成 $R$-模.将 $M$ 分解为不变量形式：
\[
  M \cong R^d \oplus R/(a_1) \oplus \dots \oplus R/(a_k) \quad (a_1 \mid \dots \mid a_k),
\]
则自由秩 $d$ 恰好等于 $\dim_K(K \otimes_R M)$（其中 $K$ 为 $R$ 的分式域）.因此自由秩 $d$ 由 $M$ 唯一定义.
\end{theorem}

\begin{theorem}[基变换保持矩阵表示]\label{thm:base-ext-matrix}
设 $M$ 与 $M'$ 为非零有限自由 $R$-模，$\varphi \colon M \to M'$ 为 $R$-线性映射.任取基 $\{e_j\}$ 与 $\{e'_i\}$，则 $S$-线性映射 $\varphi_S \colon S \otimes_R M \to S \otimes_R M'$ 关于基 $\{1 \otimes e_j\}$ 与 $\{1 \otimes e'_i\}$ 的矩阵表示，与 $\varphi$ 关于基 $\{e_j\}$ 与 $\{e'_i\}$ 的矩阵表示完全相同！
\end{theorem}

\begin{proof}
若 $\varphi(e_j) = \sum_i a_{ij} e'_i$，则
\[
  \varphi_S(1 \otimes e_j) = 1 \otimes \varphi(e_j) = \sum_i a_{ij} (1 \otimes e'_i),
\]
矩阵系数全然未变.
\end{proof}

\begin{theorem}\label{thm:s-tensor-hom}
设 $M$ 与 $M'$ 为 $R$-模.存在唯一的 $S$-线性映射
\[
  S \otimes_R \Hom_R(M, M') \to \Hom_S(S \otimes_R M, S \otimes_R M'), \quad s \otimes \varphi \mapsto s\varphi_S.
\]
当 $M$ 与 $M'$ 为有限自由 $R$-模时，该映射是 $S$-模同构.特别地，取 $M' = R$，有自然同构
\[
  S \otimes_R (M^{\vee_R}) \cong (S \otimes_R M)^{\vee_S}.
\]
\end{theorem}

\begin{theorem}[整环上自由子模单射张量幂保持定理]\label{thm:domain-free-tensor-power-injective}
设 $R$ 为整环，$\varphi \colon M \hookrightarrow N$ 为单射线性映射.若 $M$ 是自由模，则对任意 $k \geq 1$，张量幂映射
\[
  \varphi^{\otimes k} \colon M^{\otimes k} \to N^{\otimes k}
\]
必为单射.
\end{theorem}

\begin{proof}
考虑交换图：
\begin{center}
\begin{tikzcd}[row sep=large, column sep=large]
M^{\otimes k} \arrow[r, "\varphi^{\otimes k}"] \arrow[d, hook] & N^{\otimes k} \arrow[d] \\
K \otimes_R M^{\otimes k} \arrow[r, "1 \otimes \varphi^{\otimes k}"'] \arrow[d, "\cong"'] & K \otimes_R N^{\otimes k} \arrow[d, "\cong"] \\
(K \otimes_R M)^{\otimes_K k} \arrow[r, "\varphi_K^{\otimes k}"'] & (K \otimes_R N)^{\otimes_K k}
\end{tikzcd}
\end{center}
由分式域 $K$ 的平坦性，$\varphi_K \colon K \otimes_R M \to K \otimes_R N$ 是向量空间的单射，故其张量幂 $\varphi_K^{\otimes k}$ 必单.又因 $M^{\otimes k}$ 是自由模（无扭），左侧坚线复合是单射.由交换图追逐即知顶层 $\varphi^{\otimes k}$ 必为单射.
\end{proof}

% sections2/sec5_vector_spaces.tex
\section{向量空间}
\label{sec:vector-spaces}

因为所有（非零）向量空间都拥有基，我们在模论中所讨论的一般结果在向量空间背景下变得极为清晰和纯粹.

设 $K$ 为域.域上非零向量空间的张量积决不会发生意外的坍缩：若 $V$ 与 $W$ 为非零 $K$-向量空间，则 $V \otimes_K W \neq 0$，且在基数意义下恒有
\[
  \dim_K(V \otimes_K W) = (\dim_K V)(\dim_K W).
\]

对于任意 $K$-线性映射 $V \xrightarrow{\varphi} V'$ 与 $W \xrightarrow{\psi} W'$，映射 $\varphi \otimes \psi$ 保持同构、满射性，并且由于所有向量空间皆为自由模，\textbf{单射线性映射的张量积恒为单射}（定理~\ref{thm:free-is-flat}）.

\begin{example}\label{ex:subspace-tensor}
子空间的张量积就是张量积的子空间：若 $V \subset V'$ 且 $W \subset W'$，则典范线性映射 $V \otimes_K W \hookrightarrow V' \otimes_K W'$ 是单射.
\end{example}

\begin{theorem}\label{thm:subspace-tensor-equality}
设 $V \subset V'$ 且 $W \subset W'$，其中 $V'$ 与 $W'$ 非零.则 $V \otimes_K W = V' \otimes_K W'$ 当且仅当 $V = V'$ 且 $W = W'$.
\end{theorem}

\begin{proof}
由对称性，若 $V \subsetneq V'$，存在线性泛函 $\varphi \in V'^\vee$ 使得 $\varphi|_V = 0$ 但 $\varphi(v'_0) = 1$（某 $v'_0 \in V'$）.取非零 $\psi \in W'^\vee$ 使得 $\psi(w'_0) = 1$.则线性泛函 $v' \otimes w' \mapsto \varphi(v')\psi(w')$ 在 $V \otimes_K W'$ 上恒为零，但在 $v'_0 \otimes w'_0$ 处取值为 $1$.故 $v'_0 \otimes w'_0 \notin V \otimes_K W'$，说明 $V \otimes_K W' \subsetneq V' \otimes_K W'$.
\end{proof}

当 $V$ 与 $W$ 为有限维空间时，定理~\ref{thm:hom-tensor-hom} 给出典范同构：
\[
  \Hom_K(V, V') \otimes_K \Hom_K(W, W') \cong \Hom_K(V \otimes_K W, V' \otimes_K W').
\]
取 $V' = W' = K$，即得对偶空间的张量积与张量积对偶的自然同构：
\begin{equation}
  V^\vee \otimes_K W^\vee \cong (V \otimes_K W)^\vee, \quad (\varphi \otimes \psi)(v \otimes w) = \varphi(v)\psi(w).
\end{equation}
有限维时归纳可得 $(V^\vee)^{\otimes k} \cong (V^{\otimes k})^\vee$.

\begin{theorem}[线性映射张量积的核与像]\label{thm:image-kernel-vector-space}
设 $V_1 \xrightarrow{\varphi_1} W_1$ 与 $V_2 \xrightarrow{\varphi_2} W_2$ 为线性映射.则
\[
  \ker(\varphi_1 \otimes \varphi_2) = \ker \varphi_1 \otimes_K V_2 + V_1 \otimes_K \ker \varphi_2, \quad \operatorname{im}(\varphi_1 \otimes \varphi_2) = \varphi_1(V_1) \otimes_K \varphi_2(V_2).
\]
特别地，当 $V_1, V_2 \neq 0$ 时，$\varphi_1 \otimes \varphi_2$ 单当且仅当 $\varphi_1$ 与 $\varphi_2$ 均单；当 $W_1, W_2 \neq 0$ 时，$\varphi_1 \otimes \varphi_2$ 满当且仅当 $\varphi_1$ 与 $\varphi_2$ 均满.
\end{theorem}

\begin{corollary}\label{cor:subspace-quotient-tensor}
设 $V \subset V'$ 且 $W \subset W'$.则
\[
  (V' \otimes_K W') / (V \otimes_K W' + V' \otimes_K W) \cong (V'/V) \otimes_K (W'/W).
\]
\end{corollary}

\begin{theorem}\label{thm:elementary-tensor-zero}
设 $v_i \in V_i$.则在 $V_1 \otimes_K \dots \otimes_K V_k$ 中 $v_1 \otimes \dots \otimes v_k = 0$ 当且仅当至少有一个 $v_i = 0$.
\end{theorem}

\begin{proof}
若每个 $v_i \neq 0$，由线性无关性延拓可知存在 $\varphi_i \in V_i^\vee$ 使得 $\varphi_i(v_i) = 1$.则线性函数 $\varphi_1 \otimes \dots \otimes \varphi_k$ 在 $v_1 \otimes \dots \otimes v_k$ 上取值为 $1 \neq 0$，故该张量非零.
\end{proof}

\begin{theorem}[简单张量相等的充分必要条件]\label{thm:elementary-tensor-equality}
设 $v_i, v'_i \in V_i$ 均非零.则在 $V_1 \otimes_K \dots \otimes_K V_k$ 中
\[
  v_1 \otimes \dots \otimes v_k = v'_1 \otimes \dots \otimes v'_k
\]
当且仅当存在非零标量 $c_1, \dots, c_k \in K$，满足
\[
  v_i = c_i v'_i \quad (1 \leq i \leq k) \quad\text{且}\quad c_1 c_2 \dots c_k = 1.
\]
\end{theorem}

定理~\ref{thm:elementary-tensor-equality} 具有极具启发性的几何与物理涵义：简单张量 $v_1 \otimes \dots \otimes v_k$ 虽然不能唯一确定各个向量 $v_i$ 本身，但它\textbf{唯一确定了每一个一维子空间（射线方向）$K v_i$}！

在量子力学中，系统的纯态由复 Hilbert 空间中的一维子空间所描述.当组合两个系统时，复合态空间为张量积 $H_1 \otimes_\C H_2$.由简单张量生成的子空间 $\C(v_1 \otimes v_2)$ 唯一分解为各子系统的状态空间 $\C v_1$ 与 $\C v_2$（非纠缠态）；而绝大多数非简单张量所张成的状态则无法拆解为各自独立的子系统状态——这正是量子力学中最本质的核心特征：\newterm{quantum-entanglement}{量子纠缠}{quantum entanglement}（Schr\"odinger \cite{schrodinger}）.

\begin{theorem}[对称群在张量幂上的作用]\label{thm:sym-group-action}
对置换 $\sigma \in S_k$，存在唯一的线性映射 $P_\sigma \colon V^{\otimes k} \to V^{\otimes k}$，满足
\[
  P_\sigma(v_1 \otimes \dots \otimes v_k) = v_{\sigma^{-1}(1)} \otimes \dots \otimes v_{\sigma^{-1}(k)}.
\]
映射满足 $P_\sigma \circ P_\tau = P_{\sigma\tau}$；当 $\dim_K V > 1$ 时，$\sigma \mapsto P_\sigma$ 是群 $S_k$ 到 $\mathrm{GL}(V^{\otimes k})$ 的忠实表示.
\end{theorem}

群 $S_k$ 与一般线性群 $\mathrm{GL}(V)$ 在 $V^{\otimes k}$ 上的作用相互交换：$P_\sigma(g^{\otimes k}(t)) = g^{\otimes k}(P_\sigma(t))$.这一对易关系构成了李群与表示论中著名的\newterm{schur-weyl}{Schur--Weyl 对偶}{Schur-Weyl duality}.满足 $P_\sigma(t) = t$ 的张量称为\termcn{对称张量}；满足 $P_\sigma(t) = (\operatorname{sgn}\sigma)t$ 的张量称为\termcn{反对称张量}.

% sections2/sec6_contraction.tex
\section{张量缩并}
\label{sec:contraction}

设 $V$ 为域 $K$ 上的有限维向量空间.求值配对
\[
  V \times V^\vee \to K, \quad (v, \varphi) \mapsto \varphi(v)
\]
是 $K$-双线性的，因此诱导出唯一的 $K$-线性映射
\begin{equation}
  c \colon V \otimes_K V^\vee \to K, \quad c(v \otimes \varphi) = \varphi(v).
\end{equation}
该映射称为\newterm{tensor-contraction}{张量缩并}{tensor contraction}.

在基底与对偶基底下的表现如下：设 $\{e_1, \dots, e_n\}$ 为 $V$ 的一组基，$\{e^1, \dots, e^n\}$ 为 $V^\vee$ 的对偶基.若 $v = \sum v^i e_i$ 且 $w = \sum w_j e^j$，则
\begin{equation}
  c(v \otimes w) = \sum_{i, j} v^i w_j e^j(e_i) = \sum_{i} v^i w_i. \label{eq:6.1}
\end{equation}
在同构 $V \otimes_K V^\vee \cong \End_K(V)$ 下，缩并运算 $c$ 恰好对应于矩阵的\textbf{迹}（Trace）：
\[
  \Tr(T) = \sum_{i} T^i_i.
\]

更一般地，对于混合张量空间 $V^{\otimes k} \otimes (V^\vee)^{\otimes \ell}$（$k, \ell \geq 1$），选取第 $r$ 个逆变分量与第 $s$ 个协变分量进行对偶求值，诱导出一般张量缩并映射：
\begin{equation}
  c_{r, s} \colon V^{\otimes k} \otimes (V^\vee)^{\otimes \ell} \to V^{\otimes (k-1)} \otimes (V^\vee)^{\otimes (\ell-1)},
\end{equation}
在 Einstein 求和约定下，缩并简记为将一对上下重复的相同指标进行求和，例如：
\[
  T^{i_1 \dots i_k}_{j_1 \dots j_\ell} \mapsto T^{i_1 \dots m \dots i_k}_{j_1 \dots m \dots j_\ell}.
\]

当空间 $V$ 上赋予了非退化对称双线性型（度规张量）$g \colon V \xrightarrow{\cong} V^\vee$ 时，我们可以利用 $g$ 实现\termcn{降指标}（$v_j = \sum_i g_{ij} v^i$），或利用其逆 $g^{-1}$ 实现\termcn{升指标}（$v^j = \sum_i g^{ij} v_i$），从而将同类张量（如 $V \otimes_K V$）通过复合映射
\begin{equation}
  V \otimes_K V \xrightarrow{\id \otimes g} V \otimes_K V^\vee \xrightarrow{c} K
\end{equation}
进行度规缩并：$T^{ij} \mapsto \sum_{i, j} g_{ij} T^{ij}$.这是黎曼几何与广义相对论中所有标量曲率及场方程运算的数学基石.

% sections2/sec7_algebras.tex
\section{\texorpdfstring{$R$}{R}-代数的张量积}
\label{sec:algebras}

在前面的章节中，我们多次见到了代数（环）之间的张量积同构，例如 $\C \otimes_\R \mathrm{M}_n(\R) \cong \mathrm{M}_n(\C)$ 作为复向量空间同构.现在我们将证明：两个环的张量积自身可以赋予自然的环乘法结构，从而使这些同构直接在代数（环）层面上成立.

为了将矩阵代数等非交换结构纳入其中，我们允许所讨论的 $R$-代数非交换，但基础标量环 $R$ 始终假定为交换环.

一个 \newterm{r-algebra-def}{$R$-代数}{$R$-algebra} 是一个赋予了 $R$-双线性乘法 $A \times A \to A$ 的 $R$-模 $A$，满足结合律、存在乘法幺元 $1_A$，并且标量乘法与环乘法满足：
\begin{equation}
  r(ab) = (ra)b = a(rb) \quad (\forall r \in R, a, b \in A). \label{eq:7.1}
\end{equation}
由张量积的普遍映射性质，给出双线性乘法 $A \times A \to A$ 等价于给出 $R$-线性映射 $m \colon A \otimes_R A \to A$.

\begin{theorem}[代数张量积的乘法结构]\label{thm:algebra-tensor-mult}
设 $A$ 与 $B$ 为 $R$-代数.在加法群 $A \otimes_R B$ 上存在唯一的环乘法，使其成为具有幺元 $1 \otimes 1$ 的 $R$-代数，且对所有简单张量满足：
\begin{equation}
  (a \otimes b)(a' \otimes b') = aa' \otimes bb'. \label{eq:7.2}
\end{equation}
\end{theorem}

\begin{proof}
利用 $A$ 与 $B$ 自身的乘法线性映射 $m_A \colon A \otimes_R A \to A$ 与 $m_B \colon B \otimes_R B \to B$，形成张量积映射 $m_A \otimes m_B \colon (A \otimes_R A) \otimes_R (B \otimes_R B) \to A \otimes_R B$.结合张量积的结合律与交换律同构：
\[
  (A \otimes_R B) \otimes_R (A \otimes_R B) \cong (A \otimes_R A) \otimes_R (B \otimes_R B) \xrightarrow{m_A \otimes m_B} A \otimes_R B,
\]
它在简单张量上正好给出 $(a \otimes b) \otimes (a' \otimes b') \mapsto aa' \otimes bb'$.由此诱导出的乘法满足结合律且以 $1 \otimes 1$ 为幺元.
\end{proof}

\begin{corollary}\label{cor:comm-algebras-tensor}
若 $A$ 与 $B$ 为交换 $R$-代数，则 $A \otimes_R B$ 亦为交换 $R$-代数.
\end{corollary}

\begin{example}\label{ex:c-tensor-c}
考虑实代数 $\C \otimes_\R \C$.它作为实向量空间是 $4$ 维的，具有基 $\{1 \otimes 1, 1 \otimes i, i \otimes 1, i \otimes i\}$.令 $x = i \otimes 1$ 与 $y = 1 \otimes i$，则 $x^2 = -1, y^2 = -1$ 且 $xy = yx$.由于方程 $T^2 + 1 = 0$ 在其中有四个根 $(\pm x, \pm y)$，该交换环绝不是域！事实上，存在显式代数同构：
\[
  \C \otimes_\R \C \cong \C \times \C, \quad z \otimes w \mapsto (zw, z\bar{w}).
\]
\end{example}

\begin{example}\label{ex:r-tensor-q-r-zero-divisors}
代数 $\R \otimes_\Q \R$ 拥有大量的零因子：
\[
  (\sqrt{2} \otimes 1 + 1 \otimes \sqrt{2})(\sqrt{2} \otimes 1 - 1 \otimes \sqrt{2}) = 2 \otimes 1 - 1 \otimes 2 = 0,
\]
而左右两因子均非零.
\end{example}

\begin{example}\label{ex:poly-tensor-algebra}
作为 $R$-代数，有自然同构
\[
  R[X] \otimes_R R[Y] \cong R[X, Y], \quad R[X]^{\otimes k} \cong R[X_1, \dots, X_k].
\]
\end{example}

\begin{example}\label{ex:rational-function-field-tensor}
对于域 $F$ 上纯超越扩域的张量积：$F(x) \otimes_F F(y)$ \textbf{绝不是域}，因而绝不同构于 $F(x, y)$！事实上，它同构于分母可分解为 $b(x)d(y)$ 形式的有理函数子环（即局部化环）.元素 $1 \otimes 1 - x \otimes y$ 在其中不可逆（因为 $1/(1-xy)$ 无法分拆分母）.
\end{example}

\begin{theorem}[代数张量积的普遍映射性质]\label{thm:algebra-tensor-ump}
设 $A$ 与 $B$ 为 $R$-代数.对于任意 $R$-代数 $C$ 以及任意满足像元对易（即 $\varphi(a)\psi(b) = \psi(b)\varphi(a)$）的 $R$-代数同态 $\varphi \colon A \to C$ 与 $\psi \colon B \to C$，存在唯一的 $R$-代数同态
\[
  \Phi \colon A \otimes_R B \to C,
\]
满足 $\Phi(a \otimes b) = \varphi(a)\psi(b)$.
\end{theorem}

\begin{theorem}[多项式代数商环的基变换]\label{thm:poly-quotient-algebra-base-ext}
设 $f \colon R \to S$ 为交换环同态，$h(X) \in R[X]$.则作为 $S$-代数有
\begin{equation}
  S \otimes_R (R[X]/(h(X))) \cong S[X]/(h^f(X)),
\end{equation}
其中 $h^f(X)$ 表示将 $h(X)$ 的系数通过 $f$ 映到 $S$ 中所得的多项式.
\end{theorem}

\begin{example}\label{ex:field-extension-tensors}
利用定理~\ref{thm:poly-quotient-algebra-base-ext}，可以极其简便地计算数域扩张的张量积：
\begin{enumerate}
  \item $\R \otimes_\Q \Q(\sqrt{2}) \cong \R[X]/(X^2 - 2) \cong \R \times \R$；
  \item $\R \otimes_\Q \Q(\sqrt[3]{2}) \cong \R[X]/(X^3 - 2) \cong \R \times \C$；
  \item $\Q(\sqrt{2}) \otimes_\Q \Q(\sqrt{3}) \cong \Q(\sqrt{2})[X]/(X^2 - 3) \cong \Q(\sqrt{2}, \sqrt{3})$（这是一个域！）.
\end{enumerate}
\end{example}

\begin{theorem}\label{thm:intermediate-field-tensor}
设 $L/K$ 与 $F/K$ 为域扩张，且 $K \subset E \subset L$ 为中间域.则有典范环同构
\[
  F \otimes_K L \cong (F \otimes_K E) \otimes_E L.
\]
\end{theorem}

% sections2/sec8_tensor_algebra.tex
\section{\texorpdfstring{$R$}{R}-模的张量代数}
\label{sec:tensor-algebra}

一般的 $R$-模上并没有乘法运算.然而，存在一种典范构造，能够以``最自由、最通用''的方式将任意 $R$-模 $M$ 嵌入为一个 $R$-代数的生成元集.这个代数就是 $M$ 的\newterm{tensor-algebra}{张量代数}{tensor algebra} $T(M)$.

设 $M$ 为 $R$-模.对于 $m_1, m_2 \in M$，其最一般、满足双线性的形式积应当是简单张量 $m_1 \otimes m_2 \in M^{\otimes 2}$；类似地，五个元素的积应当属于 $M^{\otimes 5}$.这启发我们考虑所有张量幂的无限直和：
\begin{equation}
  T(M) := \bigoplus_{k \geq 0} M^{\otimes k} = R \oplus M \oplus M^{\otimes 2} \oplus M^{\otimes 3} \oplus \dots,
\end{equation}
其中约定 $M^{\otimes 0} = R$，$M^{\otimes 1} = M$.其元素为几乎处处为零的形式和 $\sum_{k \geq 0} t_k$（$t_k \in M^{\otimes k}$）.

\begin{theorem}[分次代数构造定理]\label{thm:graded-algebra-construction}
设 $\{M_k\}_{k \geq 0}$ 是一列 $R$-模，其中 $M_0 = R$.若存在一族 $R$-双线性映射 $\mu_{k, \ell} \colon M_k \times M_\ell \to M_{k+\ell}$，满足：
\begin{enumerate}
  \item 与 $R$ 的标量乘法相容：$\mu_{k, 0}(x, r) = rx$ 且 $\mu_{0, \ell}(r, y) = ry$；
  \item 满足结合律：$\mu_{k, \ell+n}(x, \mu_{\ell, n}(y, z)) = \mu_{k+\ell, n}(\mu_{k, \ell}(x, y), z)$；
\end{enumerate}
则直和 $\bigoplus_{k \geq 0} M_k$ 在乘法规则
\begin{equation}
  \left(\sum_{k \geq 0} m_k\right) \cdot \left(\sum_{\ell \geq 0} m'_\ell\right) = \sum_{n \geq 0} \left(\sum_{k+\ell=n} \mu_{k, \ell}(m_k, m'_\ell)\right)
\end{equation}
下构成一个具有幺元 $1 \in M_0$ 的结合 $R$-代数（称为\newterm{graded-algebra}{分次代数}{graded algebra}）.
\end{theorem}

\begin{lemma}\label{lem:tensor-power-bilinear}
对任意 $k, \ell \geq 0$，存在唯一的双线性映射 $\beta_{k, \ell} \colon M^{\otimes k} \times M^{\otimes \ell} \to M^{\otimes (k+\ell)}$，在简单张量上满足
\[
  (m_1 \otimes \dots \otimes m_k, m'_1 \otimes \dots \otimes m'_\ell) \mapsto m_1 \otimes \dots \otimes m_k \otimes m'_1 \otimes \dots \otimes m'_\ell.
\]
\end{lemma}

\begin{definition}\label{def:tensor-algebra}
由定理~\ref{thm:graded-algebra-construction} 与引理~\ref{lem:tensor-power-bilinear} 赋予乘法结构的分次代数
\[
  T(M) = \bigoplus_{k \geq 0} M^{\otimes k}
\]
称为 $R$-模 $M$ 的\termcn{张量代数}.
\end{definition}

\begin{example}\label{ex:free-tensor-algebra}
\begin{enumerate}
  \item 若 $M = R$，则 $M^{\otimes k} \cong R$，其张量代数同构于单变元多项式代数：$T(R) \cong R[X]$；
  \item 若 $M$ 是以 $e_1, \dots, e_n$ 为基的有限自由模，则 $T(M)$ 是 $n$ 个\textbf{非交换自由不定元}的多项式代数 $R\langle e_1, \dots, e_n\rangle$（因为当 $i \neq j$ 时，$e_i \otimes e_j \neq e_j \otimes e_i$）.
\end{enumerate}
\end{example}

通过规范嵌入 $i \colon M \hookrightarrow T(M)$（映至 $k = 1$ 的次数 $1$ 次分量），可将 $M$ 自然视作 $T(M)$ 的子模.张量代数满足以下核心普遍性质：

\begin{theorem}[张量代数的普遍映射性质]\label{thm:tensor-algebra-ump}
设 $M$ 为 $R$-模.对于任意结合 $R$-代数 $A$ 以及任意 $R$-线性映射 $f \colon M \to A$，存在唯一的 $R$-代数同态 $F \colon T(M) \to A$ 使得如下交换图成立：
\begin{equation}
\begin{tikzcd}[row sep=large, column sep=large]
M \arrow[r, "i", hook] \arrow[dr, "f"'] & T(M) \arrow[d, "\exists!\,F", dashed] \\
& A
\end{tikzcd}
\end{equation}
其在简单张量上的显式表达式为：
\[
  F(m_1 \otimes \dots \otimes m_k) = f(m_1) f(m_2) \dots f(m_k).
\]
\end{theorem}

定理~\ref{thm:tensor-algebra-ump} 表明：\textbf{从模 $M$ 到任意代数的线性映射，都能唯一地提升为从张量代数 $T(M)$ 出发的代数同态}.在范畴论中，这表明张量代数函子 $T(-)$ 是从代数范畴到模范畴的遗忘函子的左伴随函子.

张量代数是现代多重线性代数与几何学的基础构造源泉：
\begin{itemize}
  \item 将 $T(M)$ 对由所有形如 $m \otimes m$ 生成的双边理想作商，即得到模 $M$ 的\newterm{exterior-algebra}{外代数}{exterior algebra} $\bigwedge(M)$；
  \item 将 $T(M)$ 对由所有形如 $m_1 \otimes m_2 - m_2 \otimes m_1$ 生成的双边理想作商，即得到模 $M$ 的\newterm{symmetric-algebra}{对称代数}{symmetric algebra} $\operatorname{Sym}(M)$；
  \item 将切空间张量代数对 Clifford 关系作商，即生成\newterm{clifford-algebra}{Clifford 代数}{Clifford algebra}；对李代数的对易子关系作商，即构造出著名的\newterm{universal-enveloping}{万有包络代数}{universal enveloping algebra} $U(\mathfrak{g})$.
\end{itemize}
所有这些近代几何与现代物理学不可或缺的代数系统，均根植于张量积与张量代数这一坚不可摧的坚实基石之上.

% sections2/bibliography.tex
\begin{thebibliography}{99}
\addcontentsline{toc}{section}{参考文献}

\bibitem{friedl}
S.~Friedl, \textit{Eta invariants as sliceness obstructions and their relation to Casson--Gordon invariants}, Algebr. Geom. Topol. \textbf{4} (2004), 893--934. 
\url{https://arxiv.org/abs/math/0305402}.

\bibitem{henderson-pukelsheim}
H.~V. Henderson, F.~Pukelsheim, \textit{On the History of the Kronecker Product}, Linear and Multilinear Algebra \textbf{14} (1983), 113--120.

\bibitem{jacobson}
N.~Jacobson, \textit{Basic Algebra II}, 2nd ed., W. H. Freeman \& Co., New York, 1989.

\bibitem{schrodinger}
E.~Schr\"odinger, \textit{Discussion of Probability Relations between Separated Systems}, Math. Proc. Camb. Philos. Society \textbf{31} (1935), 555--563.

\end{thebibliography}


\printterminology

\end{document}
